% Awesome Source CV LaTeX Template
%
% This template has been downloaded from:
% https://github.com/darwiin/awesome-neue-latex-cv
%
% Author:
% Christophe Roger
%
% Template license:
% CC BY-SA 4.0 (https://creativecommons.org/licenses/by-sa/4.0/)
%Section: Work Experience at the top
\sectionTitle{EXPÉRIENCE PROFESSIONNELLE}{\faSuitcase}
\begin{experiences}
  \experience
    {À ce jour}   {Ingénieur R\&D Senior}{SCHNEIDER ELECTRIC}{Pacy-sur-Eure, France}
    {Décembre 2022} {
      \begin{itemize}
        \item Pilotage Agile et coaching d'équipes pluridisciplinaires
        \item Modélisation MBD, optimisation multi-critères de composants
        \item Génération de code embarqué via Embedded Coder
        \item Validation et intégration par tests MIL/SIL/HIL
        \item Conception de pipelines de données complexes et modèles prédictifs (ML)
        \item DevOps : CI/CD, automatisation tests, infrastructure cloud
        \item Collaboration transverse pour la conformité produit
        \item Partenariat avec équipes R\&D globales et transfert technologique
        \item Rédaction de livrables techniques et diffusion des standards d'ingénierie
      \end{itemize}
    }
    {JIRA, MATLAB, Simulink, Python, Scrum, DevOps, Hydraulique, Thermodynamique, MBD, Simulation, ML, AI}
  \emptySeparator

  \experience
    {Novembre 2022}   {Ingénieur R\&D Senior}{TRISKELL-CONSULTING}{Puteaux, France}
    {Février 2022} {
      \begin{itemize}
        \item Modélisation multiphysique (MBD) de systèmes complexes
        \item Optimisation du dimensionnement et du contrôle
        \item Conception de modèles prédictifs et traitement de données
        \item Mise en oeuvre de pratiques DevOps : pipelines CI/CD, kubernetes
        \item Documentation technique
      \end{itemize}
    }
    {MATLAB, Simulink, Python, ML, IA, DevOps, MBD, système, Agile, Stateflow, Embedded Coder, TargetLink}
  \emptySeparator

  \experience
    {Janvier 2022}   {Ingénieur R\&D Senior}{CARMAT}{Vélizy-Villacoublay, France}
    {Décembre 2021} {
      \begin{itemize}
        \item Définition d'indicateurs techniques et physiologiques
        \item Conception et maintenance de pipelines de données complexes
        \item Conception de modèles prédictifs avec Machine Learning
        \item Développement de tableaux de bord dynamiques (Power BI, Tableau)
        \item Pilotage Agile des projets et industrialisation
      \end{itemize}
    }
    {Python, Power BI, Tableau, ML, IA, DevOps, SQL, ETL, Analytics}
  \emptySeparator

  \experience
    {Novembre 2021}   {Ingénieur R\&D}{PSG-DOVER}{Auxerre, France}
    {Janvier 2019} {
      \begin{itemize}
        \item Support technique produit et mise en production
        \item Conception et maintenance de pipelines de données complexes
        \item Conception de modèles prédictifs et traitement de données
        \item Modélisation multiphysique et simulation avancée
        \item Rédaction technique et veille sectorielle
        \item Optimisation des performances
        \item Capitalisation des retours d'expérience pour fiabiliser les solutions
      \end{itemize}
    }
    {MATLAB, Simulink, Python, CATIA v5, ML, IA, Hydraulique, Thermodynamique, MBD, Simulation}
  \emptySeparator

  \experience
    {Décembre 2018}   {Ingénieur R\&D}{LUSAC}{Caen, France}
    {Décembre 2015} {
      \begin{itemize}
        \item Conception et validation expérimentale
        \item Suivi de projet et coordination des prestataires
        \item Conception de modèles prédictifs et modélisation multiphysique
        \item Développement de briques technologiques
        \item Veille technologique sur les systèmes énergétiques
        \item Présentation des résultats techniques et recommandations d'amélioration
      \end{itemize}
    }
    {MATLAB, Simulink, Python, CATIA v5, TensorFlow, Thermodynamique, Hydrogène}
  \emptySeparator

  \experience
    {Septembre 2015}   {Ingénieur R\&D}{LIEBHERR AEROSPACE AND TRANSPORTATION}{Toulouse, France}
    {Avril 2015} {
      \begin{itemize}
        \item Conception et validation expérimentale
        \item Définition des scénarios de test
        \item Conception de modèles prédictifs
      \end{itemize}
    }
    {VBA, Simulink, CATIA v5, Hydraulique, Thermodynamique, Transfert thermique, Model-Based Design, Simulation}
  \emptySeparator

  \experience
    {Mars 2015}   {Ingénieur R\&D}{ARIANESPACE}{Les Mureaux, France}
    {Décembre 2014} {
      \begin{itemize}
        \item Modélisation physique et analyse de trajectoire
        \item Conception de modèles prédictifs
        \item Modélisation multiphysique et simulation
      \end{itemize}
    }
    {Fortran, Thermodynamique, Programmation, Aérodynamique, Simulation numérique}
  \emptySeparator

  \experience
    {Juillet 2014}   {Ingénieur R\&D}{APERAM}{Gueugnon, France}
    {Février 2014} {
      \begin{itemize}
        \item Conception et dimensionnement de protecteurs industriels
        \item Réalisation de calculs et analyses pour la conformité
        \item Modélisation multiphysique et simulation
      \end{itemize}
    }
    {VBA, Simulink, CATIA v5, Thermodynamique, Structures, CAO}
  \emptySeparator
\end{experiences}
